\documentclass{article}
\usepackage{qdraw}
\usepackage{siunitx}

\DeclareSIUnit{\point}{pt}

\def\tempcode{}
\NewDocumentCommand{\codeexample}{ +v }{%
    \noindent
    \begin{tabular}{ @{} p{0.25\textwidth} @{} p{0.75\textwidth} @{} }
        \raisebox{-1.125\height}{%
            \scantokens{#1}%
        } & \texttt{#1} \\
    \end{tabular}
    \bigskip
}

\begin{document}

\section*{Cheatsheet for \texttt{qdraw}}

This cheatsheet provides an overview of the main commands and options for the \texttt{qdraw} package. It serves as a replacement for the documentation that currently does not exist yet.

\subsection*{Basic commands}

\verb|qdraw| is the main environment for drawing. The following options can be set for the environment:
\begin{itemize}
    \item \verb|baseline| Set the baseline for the environment as dimension. The default baseline is typically the bottom of the the bounding box of the drawing.
    \item \verb|layers| Set the layers for the environment as a comma-separated list. The layer \verb|main| must always be present and is the default layer. Layers are drawn in the order they are defined.
    \item \verb|overlay| Set the environment to have a zero-sized bounding box.
\end{itemize}

\verb|\qdrawset| is used to set options for paths. The following options can be set for the path:
\begin{itemize}
    \item \verb|stroke| Adds a stroke to the path.
    \item \verb|fill| Fills the path.
    \item \verb|clip| Uses the path as a clipping region for subsequent paths.
\end{itemize}

\codeexample{
\begin{qdraw}
    \path[stroke, clip] (30:1cm) -- (150:1cm) -- (270:1cm) ! ;
    \fill[fill color=red] (0,0) circle[radius=0.65cm] ;
\end{qdraw}
}

\begin{itemize}
    \item \verb|stroke color| Set the stroke color for the path. The default stroke color is the current (text) color.
    \item \verb|fill color| Set the fill color for the path. The default fill color is the current (text) color.
    \item \verb|stroke opacity| Set the stroke opacity for the path. The default stroke opacity is \num{1} (fully opaque).
    \item \verb|fill opacity| Set the fill opacity for the path. The default fill opacity is \num{1} (fully opaque).
\end{itemize}

\begin{itemize}
    \item \verb|line width| Set the line width for the path. The default line width is \qty{0.4}{\point}.
    \item \verb|line cap| Set the line cap for the path. The choices are \verb|butt|, \verb|round|, and \verb|rectangle|. The default line cap is \verb|butt|.
    \item \verb|line join| Set the line join for the path. The choices are \verb|miter|, \verb|round|, and \verb|bevel|. The default line join is \verb|miter|.
    \item \verb|miter limit| Set the miter limit for the path. The default miter limit is \num{10}.
    \item \verb|dash pattern| Set the dash pattern for the path as a comma-separated list of dimensions. Each item in the list is a dimension representing the length of a dash or a gap. The default dash pattern is empty (a solid line).
    \item \verb|dash phase| Set the dash phase for the path as a dimension. The default dash phase is \qty{0}{\point}.
\end{itemize}

\begin{itemize}
    \item \verb|corner arc| Set the corner arc for the path as a comma-separated list of at most two dimensions. The first applies to the corner formed by the path leading into the corner, and the second applies to the corner formed by the path leading out of the corner. If only one dimension is given, it applies to both corners. The default corner arc is \qty{0}{\point} (sharp corners).
    \item \verb|fill rule| Set the fill rule for the path. The choices are \verb|non-zero| and \verb|even-odd|. The default fill rule is \verb|non-zero|.
\end{itemize}

\begin{itemize}
    \item \verb|shift| Set the shift for the path as a tuple of dimensions. The first item in the tuple is the shift in the x-direction, and the second item is the shift in the y-direction. The default shift is (\qty{0}{\point}, \qty{0}{\point}).
    \item \verb|rotate| Set the rotation for the path as a dimension representing the angle of rotation. The default rotation is \qty{0}{\degree}.
    \item \verb|scale| Set the scale for the path as a comma-separated list of at most two floating-point numbers. The first item in the tuple is the scale in the x-direction, and the second item is the scale in the y-direction. If only one number is given, it is used for both directions. The default scale is (\num{1}, \num{1}).
    \item \verb|slant| Set the slant for the path as a comma-separated list of at most two floating-point numbers. The first item in the tuple is the slant in the x-direction, and the second item is the slant in the y-direction. If only one number is given, it is used for both directions. The default slant is (\num{0}, \num{0}).
\end{itemize}

\codeexample{
\begin{qdraw}
    \path[stroke, slant={0.25,0.5}] (0cm,0cm) rectangle[corner={(1cm,1cm)}] ;
\end{qdraw}
}

\verb|\path| is used to to draw a path inside the \verb|qdraw| environment. The path is defined by a sequence of coordinates and path construction operations. Each \verb|\path| command must end with a semicolon. The following operations can be used to construct the path:

\begin{itemize}
    \item \verb|[<options>]| Set options for the path. The options are the same as for \verb|\qdrawset|. Options can be set multiple times for the same path after any coordinate. Some options (such as color options or transformations) will apply to the entire path and the last specified value will be used.
    \item \verb|(<a>,<b>)| Define a coordinate with x and y parts. The first item in the tuple is the x-coordinate, and the second item is the y-coordinate.
    \item \verb|(<a>:<b>)| Define a coordinate with angle and radius parts. The first item in the tuple is the angle, and the second item is the radius.
    \item \verb|(<a>:<b>,<c>)| Define a coordinate with angle and radius parts. The first item in the tuple is the angle, and the second and third items are the radii.
    \item \verb|(<name>)| Reference a coordinate with a name. The coordinate is defined as the position of the node with the given name. Coordinates can be named by either using the \verb|point| or the \verb|node| command with the given name set via the \verb|name| option (see below).
\end{itemize}

\begin{itemize}
    \item \verb|-- (<coordinate>)| Draw a line to the next coordinate.
    \item \verb#|- (<coordinate>)# Draw a vertical, then horizontal line to the next coordinate.
    \item \verb#-| (<coordinate>)# Draw a horizontal, then vertical line to the next coordinate.
    \item \verb|.. (<control>) .. (<coordinate>)| Draw a curve to the next coordinate with a single control point.
    \item \verb|.. (<control 1>) & (<control 2>) .. (<coordinate>)| Draw a curve to the next coordinate with two control points. 
    \item \verb|!| Close the path by drawing a line from the last coordinate to the first coordinate.
\end{itemize}

\begin{itemize}
    \item \verb|circle[<options>]| Draw a circle with the current coordinate as center. The radius is specified as a dimension with the \verb|radius| option.
    \item \verb|ellipse[<options>]| Draw an ellipse with the current coordinate as center. The vectors are specified as tuples of dimensions with the \verb|vector a| and \verb|vector b| options.
    \item \verb|arc[<options>]| Draw an arc with the current coordinate as origin. The radii are specified as dimensions with the \verb|radius a| and \verb|radius b| options. Instead of radii, it is also possible to specify the vectors as tuples of dimensions with the \verb|vector a| and \verb|vector b| options. The angles are specified as dimensions with the \verb|start angle| and \verb|end angle| options.
    \item \verb|rectangle[<options>]| Draw a rectangle with the current coordinate as origin corner. The other corner is specified as a coordinate with the \verb|corner| option. If the \verb|relative| boolean option is set, the corner coordinate is interpreted as relative to the current coordinate instead of absolute.
    \item \verb|grid[<options>]| Draw a grid with the current coordinate as origin corner. The step is specified as a comma-separated list of at most two dimensions with the \verb|step| option. The first item in the tuple is the step in the x-direction, and the second item is the step in the y-direction. If only one dimension is given, it is used for both directions. The other corner is specified as a coordinate with the \verb|corner| option. 
\end{itemize}

\begin{itemize}
    \item \verb|point[<options>]| Define a coordinate at the current position and name it with the given name. The position of the node can be adjusted with the \verb|at| option which accepts a coordinate. The coordinate can be referenced later by using the name in parentheses as a coordinate (see above). 
    \item \verb|node[<options>] {<content>}| Place a node positioned at the current coordinate. The position of the node can be adjusted with the \verb|at| option which accepts a coordinate. The anchor is specified with the \verb|anchor| option which accepts a comma-separated list of at most two poles that define the anchor point as their intersection. Available poles are \verb|t|, \verb|b|, \verb|l|, \verb|r|, \verb|vc| and \verb|hc|. The width of the node can be set with the \verb|width| option. The content of the node is given as argument. The node can be named by setting the \verb|name| option.
\end{itemize}

\codeexample{
\begin{qdraw}
    \path (0,0) point[name=A, at={(1cm,1cm)}] ;
    \path[stroke] (0cm,0cm) -- (A) -- (1cm,2cm) ;
\end{qdraw}
}

\verb|\stroke| is a shortcut for \verb|\path [stroke]|

\verb|\fill| is a shortcut for \verb|\path [fill]|

\verb|\clip| is a shortcut for \verb|\path [clip]|

\verb|\node| is a shortcut for \verb|\path (0pt, 0pt) node|

\verb|\point| is a shortcut for \verb|\path (0pt, 0pt) point|

\subsection*{Other commands}

\verb|\map| is used to map over a comma-separated list of items or a range of numbers and execute code for each item. The code can reference the current item with \verb|#1|. The following options can be set for the command:

\begin{itemize}
    \item \verb|list| Set the command to map over a comma-separated list of items. Each item in the list is passed as an argument to the code.
    \item \verb|start| Set the start of the range for mapping over a range of numbers as a floating-point number. The default start is \num{1}.
    \item \verb|end| Set the end of the range for mapping over a range of numbers as a floating-point number. The default end is \num{1}.
    \item \verb|step| Set the step for the range for mapping over a range of numbers as a floating-point number. The default step is \num{1}.
\end{itemize}

\codeexample{
\begin{qdraw}
    \fill[fill color=red] map[list={30,150,270}] { (#1:0.75cm) circle[radius=2pt] } ;
\end{qdraw}
}

\codeexample{
\begin{qdraw}
    \map[list={30,150,270}] { [fill, fill color=red] (#1:0.75cm) circle[radius=2pt] } ;
\end{qdraw}
}

\codeexample{
\begin{qdraw}
    \path[stroke, stroke color=blue] (1cm,-1cm) map[list={30,150,270}] { [fill, fill color=red] -- (#1:0.75cm) circle[radius=2pt] } ;
\end{qdraw}
}

\end{document}